% The main text, shared by every version. Title, authors, abstract and the
% back-matter statements are in metadata.tex. Keep Methods as a single
% \section{Methods}: the Nature build moves that section to the end.

\section{Introduction}\label{sec:intro}

This is the template's example paper. Each part of it, from the sections,
equations, figure and table to the citations and references into the
supplement, is here to show a feature, and the text explains how that feature
is written and built. Replace it with your own paper.

A paper is now read in several forms: a preprint, a copy with line numbers for
co-authors and reviewers, and one or more journal submissions. Most journals
accept almost any PDF at first submission, but at acceptance they want editable
source, in practice LaTeX \citep{lamport1994} or Word. The template therefore
keeps a single LaTeX source and builds every version from it, with each
journal's formatting kept out of the text.

\section{Methods}\label{sec:methods}

\subsection{Source files}

The writing lives in three files. \texttt{metadata.tex} holds the title,
authors, affiliations, abstract and statements; \texttt{paper.tex} holds this
text; and \texttt{supplement.tex} holds the supplementary information.
References are kept in \texttt{refs.bib} and cited by key, with BibTeX
formatting them in each journal's style \citep{patashnik1988}. Figures are
produced elsewhere and copied into \texttt{figures/} as finished files.

\subsection{Building each version}

Each version has a short wrapper that sets the document class and layout,
defines how the authors in \texttt{metadata.tex} are set in that class, and
includes the shared text. Running \texttt{make} compiles each wrapper with
pdfLaTeX and BibTeX (Fig.~\ref{fig:pipeline}). For a journal it can also make
a submission bundle: the wrapper flattened into one file with the bibliography
pasted in and references into the supplement written as plain numbers,
together with the figures and the journal's class files. The Nature build
moves this Methods section after the Discussion, as Nature Portfolio journals
require, which is why Methods must stay a single section.

\subsection{Equations}

Equations use standard \texttt{amsmath} environments and are numbered in every
version. For example, the logistic model of population growth
\citep{verhulst1838} is
\begin{equation}
  \frac{\dd N}{\dd t} = r N \left(1 - \frac{N}{K}\right),
  \qquad N(0) = N_0 ,
  \label{eq:logistic}
\end{equation}
where $r$ is the growth rate and $K$ the carrying capacity. Equation~\eqref{eq:logistic}
has the solution
\begin{equation}
  N(t) = \frac{K}{1 + \left(K/N_0 - 1\right) e^{-r t}} .
  \label{eq:solution}
\end{equation}
Notation shortcuts such as the upright differential in
Eq.~\eqref{eq:logistic} are defined once in \texttt{macros.tex}.

\section{Results}\label{sec:results}

Figure~\ref{fig:pipeline} shows how the source files become the finished
versions, and Table~\ref{tab:versions} compares the example versions. Tables
use \texttt{booktabs} rules, and the grouped column headings show that
spanning columns survive every build.

\linelabel{rev:linelabels}A sentence can be marked with a line label, which
lets a response to reviewers quote its line number in the review copy. The
main text can also refer into the supplement: an example data figure is
Supplementary Fig.~\ref{S:fig:growth}, and Supplementary
Note~\ref{S:sec:crossrefs} explains how references between the two documents
are resolved.

\begin{figure}[tbp]
  \centering
  \includegraphics[width=\linewidth]{pipeline}
  \caption{\textbf{From source files to finished versions.} The shared text
    and metadata are included by one wrapper per version, which
    \texttt{make} compiles. Submission bundles are also flattened into a
    single file with the bibliography and supplementary numbers written in.
    Everything finished goes to \texttt{output/}.}
  \label{fig:pipeline}
\end{figure}

\begin{table}[tbp]
  \centering
  \caption{\textbf{How the example versions differ.} The text is identical in
    all of them.}
  \label{tab:versions}
  \begin{tabular}{lllll}
    \toprule
            & \multicolumn{2}{c}{Layout} & \multicolumn{2}{c}{Citations} \\
    \cmidrule(lr){2-3}\cmidrule(lr){4-5}
    Version & Line numbers & Figures & Style & Order \\
    \midrule
    Preprint       & No  & In text  & Numbered    & Citation \\
    Review         & Yes & In text  & Numbered    & Citation \\
    Nature         & Yes & In text  & Numbered    & Citation \\
    PLOS           & Yes & Separate & Vancouver   & Citation \\
    Bioinformatics & Yes & In text  & Author--year & Alphabetical \\
    \bottomrule
  \end{tabular}
\end{table}

\section{Discussion}\label{sec:discussion}

The template trades a few conveniences for portability. Citations are made
with one command so that numbered and author--year styles both work, and the
text avoids packages that journal classes or pandoc \citep{macfarlane2006} do
not support. It deliberately stops short of producing a finished journal
file: the journal folders are examples to be refreshed from each publisher's
current template, and the accepted version will usually need some tuning by
hand, for example replacing the few metadata commands with plain text where a
journal asks for no personal macros. The Word copy is meant for sharing drafts
rather than for submission.
